% Main Preamble

\documentclass[12pt, letterpaper]{article}
\title{Dirtywork Manual\\[0.12in] \large Version 0.1.1}
\author{Caleb Griffin}
\date{June, 2026}

\usepackage{
	xcolor, 
	graphicx, 
	titlesec, 
	geometry, 
	listings, 
	amsmath,
	hyperref
}
\usepackage[default]{lato}
\geometry{left= 0.5in, right=0.5in, top=0.5in, bottom=0.75in}

\pagecolor{white}

\hypersetup{
	colorlinks=true,
	linkcolor=black,
	filecolor=magenta,
	urlcolor=blue
}

\begin{document}
\maketitle
\thispagestyle{empty}

\pagebreak
\tableofcontents
\pagebreak

\section{Introduction}
Dirtywork is a chess engine that solves for a recommended move and an
evaluation of the game in usits of pawns advantage.

\section{User Guide}
Dirtywork can be easily run in the terminal with flags.

\subsection{Flags}
\begin{itemize}
\item -f/--file Pass a game file to dirtywork.
\item -d/--debug
\item -m/--memory Set the max memory dirtywork can request.
\end{itemize}

\subsection{TUI Guide}
Simply run dirtywork to enter the TUI.

\section{Algorithms}
The algorithm used by dirtywork is temporal minmax eval, or ``TME''.

\subsection{Depth}
The depth parameter of TME sets how many times the algorithm is recursivly
called. A higher depth is more expensive, but more accurate. If we call
the depth parameter \textbf{n}, then we can derive the rest of the parameters
from this \textbf{n}. The prune function removes the lowest eval move
sequences and leave only the top \(2 + 2\sqrt{n}\) sequences. Prune is called
every \(\sqrt{n + 2}\) moves.

\section{Math}
Dirtywork uses bitboards to store piece locations. Bitboards ar 64 bit
integers wher each bit coresponds to a square on the board. If the bit is
zero, no peice is present, if the bit is 1, a piece is present at that square.

\subsection{Determining the area under a function}
If you want to determine \(\int_a^b x^2\,dx\) for your system.

\end{document}
